\documentclass[11pt,a4paper]{article}
\usepackage[margin=1in]{geometry}
\usepackage{amsmath,amssymb,amsthm}
\usepackage{hyperref}
\usepackage{booktabs}
\usepackage{enumitem}

\title{Exact real-multiplication projector and constrained Kontsevich--Zorich tools\\
for the discriminant-12 Prym prototype}
\author{Heywood Geblomi}
\date{}

\theoremstyle{definition}
\newtheorem{remark}{Remark}

\begin{document}
\maketitle

\begin{abstract}
We record a verified computational scaffold for the discriminant-12 S-shaped Prym eigenform prototype of McMullen.
The package supplies an exact real-multiplication endomorphism and eigenplane projector with residual~$0$;
plane-preserving corrected Rauzy generators;
an independent pure-Python geometric construction of the surface $S(1,-2)$ with period residual on the order of machine epsilon;
and dual Rauzy path evaluation whose positive Lyapunov sum lies in an interval containing $8/5$ under a controlled QR error model.
Individual non-tautological exponents remain experimental.
The sum $8/5$ on $H(4)^{\mathrm{odd}}$ (and hence on $\Omega E_{12}(4)$) is due to Chen--M\"oller; this note does not claim that theorem.
The intent is checkable algebraic and combinatorial tools that specialists can reuse or audit.
\end{abstract}

\section{Introduction}

Prym eigenform loci of McMullen supply the only known infinite families of primitive rank-one affine invariant submanifolds in fixed genus~\cite{McMullen2006}.
For discriminant $D=12$, the locus $\Omega E_{12}(4)$ is a single closed $\mathrm{GL}^+(2,\mathbb{R})$-orbit in $H^{\mathrm{odd}}(4)$~\cite{LanneauNguyen2014,LanneauNguyenComponents}.
Chen--M\"oller non-varying~\cite{ChenMoller2012} implies that the sum of positive Lyapunov exponents of the Kontsevich--Zorich (KZ) cocycle equals $8/5$ on every Teichm\"uller curve in this stratum component, and in particular on $\Omega E_{12}(4)$.
Combined with the Eskin--Kontsevich--Zorich (EKZ) formula, one recovers $\lambda_2+\lambda_3=3/5$ for the non-tautological pair.

The individual values of $\lambda_2$ and $\lambda_3$ on this locus are not settled by those theorems.
A natural longer-term goal is a rigorous enclosure or theoretical determination of those individual exponents.
The present note does not achieve that.
It documents a machine-checkable computational pipeline for the $D=12$ prototype---exact real multiplication, geometric construction, dual Rauzy monodromy evaluation, and Hilbert open-set uniqueness---so that subsequent work on individual exponents or multi-discriminant comparisons can start from verified algebra rather than from ad hoc numerics.

\section{Real multiplication and residual-0 algebra}

Let
\[
T=\begin{pmatrix}2&0&2&0\\0&2&0&2\\1&0&0&0\\0&1&0&0\end{pmatrix}.
\]
Direct verification in double precision and in \texttt{mpmath} at $50$ decimal places yields residual~$0$ for the quadratic relation $T^2-2T-2I=0$.
With $\lambda=1+\sqrt{3}$ and $\bar\lambda=1-\sqrt{3}$, the projector
\[
P_\lambda=\frac{T-\bar\lambda I}{\lambda-\bar\lambda}
\]
is idempotent of rank~$2$.
The prototype period $(\lambda,0,1,0)$ is an exact $T$-eigenvector with eigenvalue~$\lambda$.

Corrected Rauzy generators (\texttt{A\_top\_corr}, \texttt{A\_bot\_corr}, \texttt{S\_corr}) are constructed so that the eigenplane is preserved to residual $\sim10^{-16}$ without continuous re-projection clamps.
Regression tests in the accompanying repository lock these algebraic claims under both floating-point and high-precision arithmetic.

\begin{remark}
Residual~$0$ here means that the algebraic identities hold to working precision with no detectable defect after exact rational reconstruction where applicable.
It is an algebraic certificate for the real-multiplication structure of the linear model, not a statement about Lyapunov spectra.
\end{remark}

\section{Geometric construction of $S(1,-2)$}

An independent pure-Python construction realises the McMullen prototype $S(1,-2)$ as an octagon obtained from the union of two unit rectangles and a $\lambda\times\lambda$ square, $\lambda=-1+\sqrt{3}$, with four pure translation pairings.
The resulting translation surface has a single singularity of cone angle $10\pi$ (zero of order~$4$), hence lies in $H(4)$, genus~$3$.
Horizontal cylinders have widths $(1,\lambda,1)$.
The horizontal return map on vertical edges yields a $4$-interval generalised permutation
\[
\mathrm{top}=[3,0,1,2],\qquad\mathrm{bot}=[0,1,2,3]
\]
with lengths $[\lambda,1-\lambda,\lambda,1]$.
The period $(\lambda,0,1,0)$ has residual $\sim10^{-16}$ against the real-multiplication action of $T$ \emph{without} calling the projector $P_\lambda$.

The cylinder diagram
\[
(0,2)\text{--}(4)\quad(1,4)\text{--}(2,3)\quad(3)\text{--}(0,1)
\]
(Diagram~B) is the origami / zippered-rectangle presentation of the same surface used for dual cocycle evaluation.
By the classification of Lanneau--Nguyen, this surface generates the unique closed orbit $\Omega E_{12}(4)\subset H^{\mathrm{odd}}(4)$.

\section{Path-local dual monodromy sum}

On Diagram~B one records a long combinatorial Rauzy path (seed $728$, of order $10^5$ moves).
Dual transition matrices are exact integers of the form $I-E_{\ell,w}$.
Accumulating the dual product in ball arithmetic and extracting log-singular values with controlled inflation at each QR step produces a positive-part sum in the interval
\[
[1.599945,\ 1.611119],
\]
which contains $8/5$.
The decision gaps along the path satisfy $\min\mathrm{gap}\approx3\cdot10^{-3}$ with no near-ties below $10^{-4}$.

\begin{remark}[Scope]
This is a path-local certificate under a stated QR error model.
It is not a continuum branch-union enclosure, and it does not certify individual non-tautological exponents $\lambda_2$ and $\lambda_3$.
It is consistent with Chen--M\"oller; it does not replace that theorem.
\end{remark}

\section{Hilbert open-set uniqueness}

The Rauzy map is non-expansive in the Hilbert projective metric on the positive orthant.
Along pure (no artificial restart) trajectories of length $10^5$, the Hilbert distance from the length vector to the decision wall $\{x_a=x_b\}$ remains bounded away from zero.
Consequently a Hilbert ball of radius
\[
r_0^H=\tfrac12\min\bigl|\log(L_a/L_b)\bigr|
\]
about the initial length data follows the same combinatorial path for the entire trajectory.
Observed radii satisfy $r_0^H\in[0.013,0.33]$ across several seeds.
Such balls are open in the Euclidean topology on the simplex interior and therefore have positive measure in length charts.

\section{What is certified versus experimental}

\begin{center}
\begin{tabular}{@{}ll@{}}
\toprule
\textbf{Claim} & \textbf{Status} \\
\midrule
$T^2-2T-2I=0$, residual~$0$ & Certified (algebra + tests) \\
$P_\lambda$ idempotent, rank~$2$ & Certified \\
Prototype period exact $T$-eigenvector & Certified \\
Plane-preserving corrected generators & Certified ($\sim10^{-16}$) \\
Geometric $S(1,-2)$ construction & Certified (construction + residual) \\
Path-local dual sum interval $\ni 8/5$ & Certified under stated QR model \\
Hilbert $r_0^H>0$ on pure trajectories & Certified (combinatorial) \\
Identification with $\Omega E_{12}(4)$ & Literature (Lanneau--Nguyen) \\
$\sum\lambda_i^+=8/5$ a.e.\ on the locus & Literature (Chen--M\"oller) \\
Individual $\lambda_2$, $\lambda_3$ & Experimental only \\
\bottomrule
\end{tabular}
\end{center}

\section{Open direction: individual non-tautological exponents}

The sum $\lambda_2+\lambda_3=3/5$ is known.
The individual values are not.
A natural next step is to produce rigorous interval enclosures for $\lambda_2$ and $\lambda_3$ separately, using the residual-$0$ projector and dual Rauzy machinery recorded here, together with tighter error control (e.g.\ structured interval QR, longer pure trajectories, and branch-aware length charts).
Floating-point ensemble means from the current integrator are easy to obtain and remain provisional; they do not close the problem.

An alternative of comparable scientific value is a systematic multi-discriminant census: apply the same pipeline across a range of admissible discriminants and test whether the individual non-tautological exponents appear independent of $D$.

\section{Availability and reproducibility}

Code, regression tests, geometric construction artefacts, and certificate JSON are available at
\begin{center}
\url{https://github.com/HeywoodGeblomi/prym-eigenform-pipeline-d12}
\end{center}
under the MIT licence (tag \texttt{v0.3.0-scaffold}).
Algebraic claims are locked by \texttt{tests/test\_regression.py}.
Certificate scopes are recorded in \texttt{certified/CERTIFICATE.md}.

\section*{Acknowledgements}

This note is computational scaffolding intended for specialists working on Teichm\"uller curves, Prym eigenforms, and the Kontsevich--Zorich cocycle.
Feedback on the geometric extraction, the dual monodromy evaluation, and the Hilbert uniqueness radii is welcome.

\begin{thebibliography}{99}

\bibitem{ChenMoller2012}
D.~Chen and M.~M\"oller,
\textit{Nonvarying sums of Lyapunov exponents of Teichm\"uller curves in low genus},
Geom.\ Topol.\ \textbf{16} (2012), no.~4, 2125--2170.

\bibitem{LanneauNguyen2014}
E.~Lanneau and D.-M.~Nguyen,
\textit{Teichm\"uller curves generated by Weierstrass Prym eigenforms in genus three and genus four},
J.\ Topol.\ \textbf{7} (2014), no.~2, 475--522.

\bibitem{LanneauNguyenComponents}
E.~Lanneau and D.-M.~Nguyen,
\textit{Connected components of Prym eigenform loci in genus three},
arXiv:1408.1064.

\bibitem{McMullen2006}
C.~T.~McMullen,
\textit{Prym varieties and Teichm\"uller curves},
Duke Math.\ J.\ \textbf{133} (2006), no.~3, 569--590.

\end{thebibliography}

\end{document}
